\documentclass[literature.tex]{subfiles}

\begin{document}
\prose{Devatenáct set osmdesát čtyři}{George Orwell}
\teorieA{
Epický román, sci-fi s utopickými prvky, které ale kritizuje.
}{
Realistický román s \texthl{antiutopickým} dějem a prvky satiry.
}{
Román je psán jednoduchým a spisovným jazykem.
Speciálním prvek zde činí \enquote{novomluva},
tedy soubor speciálních slov, sloužících jako ukázka potlačování.
Orwell při scénách dbá na popis detailů.
}{}{}{
Text je pevně rozdělen na tři části.
Děj je vyprávěn er-formou se zaměřením na Winstona.
Text~je~pesimistický, čtenář jasně cítí, co cítí Winston
\textbl{$\Rightarrow$}~\texthl{útlak}, \texthl{beznaděj}, naprosté \texthl{pohlcení}.
}{
Winston viděl O’Briena možná tucetkrát během právě tolika let.
Silně ho to k němu táhlo, nejen proto, že ho zaujal rozpor mezi O’Brienovým uhlazeným chováním
a jeho zápasnickou konstitucí.
Mnohem spíš to způsobovala skrytá víra – nebo možná ani ne víra, pouze naděje –,
že O’Brienova politická uvědomělost není bez vady.
Cosi v jeho tváři to neodbytně naznačovalo.
A třeba se mu ve tváři neodrážela ani neuvědomělost, nýbrž jen inteligence.
Ale tak či onak působil dojmem člověka, se kterým se dá mluvit,
pokud se vám nějak podaří obelstít telestěnu a zastihnout ho o samotě.
Winston si tuhle domněnku nikdy ani v nejmenším nezkoušel ověřit: nebylo totiž ani možno jak.
V tu chvíli O’Brien pohlédl na hodinky, zjistil, že se blíží jedenáct nula nula,
a zjevně se rozhodl zůstat v Archivním oddělení, dokud Dvouminutovka nenávisti neskončí.
Usadil se na židli v téže řadě jako Winston, o pár míst dál.
Mezi nimi seděla drobná, písková blondýna, která pracovala v kóji vedle Winstona.
Tmavovláska se uvelebila přímo za ním.
Vzápětí se z velké telestěny na konci místnosti vydral příšerný, řezavý jek,
jako by se nějaký obludný stroj točil bez mazání.
}
\teorieB{
\wtem{Winston Smith}{hlavní postava, takzvaný every-man archetyp.
Člen \texthl{Vnější strany}.
Pracuje na \texthl{Ministerstvu Pravdy}, falšuje historické záznamy, stejně jako noviny.
Přesto ale doufá ve svobodu.
Je inteligentní, ale ne prohnaný.
Je zranitelný, ale ne slabý.
Prostě přežívá.}
\wtem{Julia}{romantický interest pro Winstona.
Prosazuje životní styl tichého odporu.
\texthl{Pragmatická} a \texthl{hedonistická}.
Kdy umí jít proti systému v malých věcech, ale nevidí smysl v organizovaném odporu.}
\wtem{O'Brian}{Winstonův \texthl{protipól}.
Winston z něj cítí, že ten jediný mu porozumí.
A on mu \texthl{skutečně rozumí}.
Ale opačně.
A chce ho dostat na svojí stranu.}
\wtem{Velký bratr}{archetyp totalitního mocnáře.
Není jasné, zda takový člověk existuje, či zda jde o další předmět moci Strany.
Jedná se o velkého a oslavovaného vůdce Strany.
Ať je to s jeho existencí jakkoli, nyní je \texthl{nesmrtelný} a \texthl{nezničitelný}.}
\wtem{Emmanuel Goldstein}{základ v reálném člověku.
Strana jeho památky využívá jako \texthl{symbol zbytečného odporu}.
Goldstein vede takzvané Bratrstvo proti Straně, ale prohrává.
A jakýkoli problém uvnitř Strany padá na jeho hlavu.
On za to může.
Ukažme si na něj.}
\wtem{Pan Charrington}{Může za finální dopadení Julie a Winstona.
Poskytl Winstonovi bezpečné útočiště.
Pocit odporu ke straně.
A byl to on, kdo celou dobu pracoval pro Stranu.
A jakékoli Winstonovi pokroky hlásil.}
}{
Příběh sleduje Winstonův běžný život a seznamuje nás s životem v Oceánii.
Jedné ze tří supervelmocí.
Možná je rok 1984, možná ne.
Informace drží Strana a ani její členi k nim nemají přístup.
První výraznou změnou je, když se Winston dá dohromady s Julií, do které se dle vlastních slov zamiloval.
Winston konečně vzdoruje Straně.
S Julií zažívá fyzické potěšení, které Strana zakazuje.
Začíná si víc věřit.
Užívá si života.
Druhou výraznou změnou je, když Winstona s Julií dostihne Psychopol a následně jsou mučeni na \texthl{Ministerstvu Lásky}.
Všechny iluze jsou rozbity.
Strana~je~všudypřítomná.
Věděla~o~každém Winstonově kroku.
Užíval~si~pouze za~jejího dohledu.
Došel jen tam, kam měl dovoleno.
Jeho odpor byl jen iluzí.
Děj končí kompletním obratem Winstonovi osobnosti.
\texthl{Kdy se kompletně podvolí Straně}.
}{
Většina příběhu se odehrává v Londýně, ve státu zvaném Oceánie.
Mohl by být rok 1984, i když to není jisté.
Během příběhu je pár vzpomínkových sekvencí a informativní pasáž v podobě Knihy, kterou Winston čte.
K tomu je příběh doprovázen Winstonovými monology, které umožňují nahlédnout do psychiky člověka pod kontrolou Strany.
}{
Autor varuje před technologií a čeho by byla schopna ve špatných rukou.
Varuje před totalitním státem a jaké kontroly by byl schopen s pomocí moderních technologií. 
}\newpage
\historie{
Autor napsal svůj román po dvou světových válkách, v období velkého technologického rozmachu.
Svět byl nestálý, v chaosu a do toho zažili největší rozkvět techniky za posledních několik staletí.
Pro románovou totalitní nadvládu ho inspiroval socialismus v SSSR.
}{}{}{}{}{
Vlastním jménem \textbl{Eric Arthur Blair}.
Dle některých zdrojů \texthl{nejvlivnější esejista své doby}.
Narodil se v chudém prostředí.
Odtud to zaměření na nejspodnější vrstvu.
Svůj pseudonym si vymyslel dle názvu řeky v místě, kde bydlel.
Svojí tvorbou se snažil upozornit na stav, ve kterém se svět nachází a mohl by se nacházet, pokud bude pokračovat směrem, kam míří.
}{
Mezi nejznámnější díla patří \textsc{Farma zvířat}, \textsc{Cesta k Wigan Pier} a \textsc{Hold Katalánsku}.
}
\kritika{}{}{
Témata jako cenzura, války za mocí, technologická nadvláda, sledování,
ztráta soukromí a další, jsou stále aktuální, ne-li ještě více, než v minulosti.
}
\nazor{
Román mistra Orwella jsem měl problém ocenit.
První dvě části byli natolik zmatečné a ploché, že můj názor byl spíše negativní.
Ale třetí část byla schopná zodpovědět každý můj problém (či na mě zapůsobila O’Brienova převýchova a já si vyvinul \texthl{dipsych})
a během posledního sta stránek se můj postoj otočil o~180~stupňů. \\
Teď ještě s delším odstupem, kdy kontroluji své spisy a opravuji chyby, kterých jsem se dopustil,
tak musím říct, že je \textbl{Devatenáct set osmdesát čtyři} vynikajícím románem.
Dopustil jsem se té chyby, že jsem psal soudy, zatímco jsem byl pln emocí z románu samotného.
A i když to nebyla chyba ve smyslu chyby, protože ve mě román vyvolal přesně ty pocity,
o kterých věřím, že je autor zamýšlel vyvolat,
tak jsem kvůli tomu hodnotil román, stejně jako pana Blaira, až moc přísně.
}
\explain{
\textbl{archetyp --} soubor vlastností vytvářející známou postavu;
\textbl{$\Rightarrow$}~postava rozpoznatelná svým chováním napříč příběhy;
\textbl{every-man --} archetyp běžného člověka potíkajícího se~s~běžnými problémy svého světa;
\textbl{pragmatismus --} klade důraz na praktické důsledky myšlení a jednání;
\textbl{hedonismus --} považuje slast za nejvyšší dobro a cíl lidského života;
\textbl{antiutopismus --} literární směr, který zobrazuje negativní vizi budoucnosti, často s totalitním režimem;
\textbl{Novomluva --} umělý jazyk vytvořený totalitním režimem s~cílem omezit myšlení a~vyjádření;
\textbl{Dipsych --} chopnost udržet současně dvě protichůdné myšlenky a věřit v obě.
}
\wpage
\end{document}